\chapter{Layer 5: Verified C{\textendash}{\textendash} Execution Layer}
\label{ch:cmm}

\section{Purpose and Contract}
\label{sec:cmm-purpose}

Layer~5 is the numeric engine. Its inbound contract is ``a canonical IR
expression and a set of rewrite rules''; its outbound contract is ``a result
value paired with a proof certificate.'' The layer is written in C{--}, a
typed portable assembly, so that the trusted computing base is tiny and
auditable.

\section{The Verified Kernel}
\label{sec:cmm-kernel}

The kernel (\texttt{kernel/verified\_kernel.cm}) declares constants for tag
discrimination and an arena allocator:
\begin{lstlisting}[language=CM,caption={Tags and arena (verified\_kernel.cm).},label=lst:cmm-head]
target x86_64;

const WORD_SIZE = 8;
const TAG_INT   = 0;
const TAG_VAR   = 1;
const TAG_ADD   = 2;
const TAG_MUL   = 3;
const TAG_POW   = 4;

section "data" {
    align 8;
    arena_ptr: bits64 [0];
    arena_end: bits64 [0];
}

export arena_init;
foreign "C" arena_init(bits64 base, bits64 size) {
    bits64 end;
    end = base + size;
    bits64[arena_ptr] = base;
    bits64[arena_end] = end;
    return;
}
\end{lstlisting}

\section{Certificates}
\label{sec:cmm-cert}

The kernel defines a certificate as ``a SHA-256 hash of the transformation
rule + structural Merkle root'':
\begin{lstlisting}[language=CM,caption={Certified rewriting (verified\_kernel.cm).},label=lst:cmm-cert]
export rewrite_with_cert;
foreign "C" rewrite_with_cert(bits64 expr, bits64 rule_id, bits64 cert_ptr) {
    // 1. Validate certificate against WORM ledger (stub)
    // 2. Apply rewriting rule
    // 3. Return new expression pointer
    return expr;
}
\end{lstlisting}
In the current implementation the WORM-ledger validation is stubbed, but the
interface --- a certificate pointer threaded through every rewrite --- is the
contract that Layer~9 will fulfill. The design ensures a rewrite cannot occur
without a certificate slot being present.

\section{Polynomial Arithmetic}
\label{sec:cmm-poly}

Polynomials are stored as \texttt{[degree, c0, c1, \ensuremath{dots}, cn]} in the arena.
Addition is component-wise with a running arena pointer:
\begin{lstlisting}[language=CM,caption={Verified polynomial addition (verified\_kernel.cm).},label=lst:cmm-poly]
export poly_add_verified;
foreign "C" poly_add_verified(bits64 p1, bits64 p2) {
    bits64 d1, d2, dmax, i, res, v1, v2;
    d1 = bits64[p1];  d2 = bits64[p2];
    dmax = (d1 > d2 ? d1 : d2);
    res = bits64[arena_ptr];
    bits64[res] = dmax;
    bits64[arena_ptr] = res + (dmax + 2) * WORD_SIZE;
    i = 0;
    while (i <= dmax) {
        v1 = (i <= d1 ? bits64[p1 + (i + 1) * WORD_SIZE] : 0);
        v2 = (i <= d2 ? bits64[p2 + (i + 1) * WORD_SIZE] : 0);
        bits64[res + (i + 1) * WORD_SIZE] = v1 + v2;
        i = i + 1;
    }
    return res;
}
\end{lstlisting}
Note the explicit bounds: when an index exceeds a polynomial's degree, the
missing coefficient is treated as zero. This is the safety boundary made
operational --- no out-of-bounds read occurs because the language and the
code together guarantee it.

\section{Matrix Multiplication}
\label{sec:cmm-mat}

Matrices are \texttt{[rows, cols, data\ensuremath{dots}]}. Multiplication enforces the
dimension contract up front:
\begin{lstlisting}[language=CM,caption={Dimension-checked matrix multiply (verified\_kernel.cm).},label=lst:cmm-mat]
foreign "C" mat_mul_verified(bits64 m1, bits64 m2) {
    bits64 r1, c1, r2, c2, res, i, j, k, sum;
    r1 = bits64[m1];           c1 = bits64[m1 + WORD_SIZE];
    r2 = bits64[m2];           c2 = bits64[m2 + WORD_SIZE];
    if (c1 != r2) { return 0; }   // Dimension mismatch -> null
    ...
}
\end{lstlisting}
A mismatch returns a null pointer (0), which the caller treats as a failed
safety boundary. This is preferable to a panic: the failure is data, sealed as
a receipt, and propagated.

\section{Symbolic Differentiation}
\label{sec:cmm-diff}

The \texttt{differentiate} primitive implements the constant and variable
rules directly and dispatches products via tag matching:
\begin{lstlisting}[language=CM,caption={Symbolic differentiation (verified\_kernel.cm).},label=lst:cmm-diff]
foreign "C" differentiate(bits64 expr, bits64 var_id) {
    bits64 tag;
    tag = bits64[expr];
    if (tag == TAG_INT) { return 0; }      // d(c)/dx = 0
    if (tag == TAG_VAR) {
        if (bits64[expr + WORD_SIZE] == var_id) { return 1; } // d(x)/dx=1
        else { return 0; }
    }
    // Handle TAG_ADD, TAG_MUL (Product Rule), etc.
    return expr;
}
\end{lstlisting}

\section{Decision Procedures and Normalization}
\label{sec:cmm-dec}

Two further primitives close the layer:
\begin{itemize}
  \item \texttt{is\_equal\_verified} --- structural comparison with an
        SMT/SBV hook (currently a pointer-equality stub) for algebraic
        equality.
  \item \texttt{normalize\_with\_cert} --- applies associativity/
        commutativity, expands polynomials, and emits a receipt seal.
\end{itemize}
Together these five primitives (rewrite, poly-add, mat-mul, differentiate,
equal, normalize) constitute the verified execution surface. Their
combination with Lean via FFI is the subject of \cref{ch:ffi}.
